\subsection{The support of a module}

DEFINITION 5. Let A be a ring and M an A-module. The set of prime ideals p of A such that \(M_{p} \neq 0\) is called the support of M and is denoted by Supp(M).

As every maximal ideal of A is prime, it follows immediately from § 3, no. 3, Corollary 2 to Theorem 1, that for A-module M to be equal to 0, it is necessary and sufficient that \(\operatorname{Supp}(\mathbf{M}) = \varnothing\) .

Example. Let a be an ideal of A; in the notation of no. 3, we have

\[
\mathbf {V} (\mathfrak {a}) = \operatorname{Supp} (\mathrm{A} / \mathfrak {a}).\tag{9}
\]

If \(p\) is a prime of \(A\) such that \(a \not\subset p\) , then \((A/a)_p = 0\) ( \(\S 3\) , no. 1, Remark 3); if on the other hand \(a \subset p\) , \(aA_p\) is contained in the maximal ideal \(pA_p\) of \(A_p\) and \((A/a)_p\) is isomorphic to \(A_p/aA_p\) and hence is non-zero ( \(\S 3\) , no. 1, Proposition 3); whence our assertion.

In particular, \(\operatorname{Supp}(\mathbf{A}) = \operatorname{Spec}(\mathbf{A})\) .

PROPOSITION 16. Let A be a ring and M an A-module.

\begin{enumerate}
\def\labelenumi{(\roman{enumi})}
\tightlist
\item
  If \(\mathbf{N}\) is a submodule of \(\mathbf{M}\) , then
\end{enumerate}

\[
\operatorname{Supp} (\mathbf {M}) = \operatorname{Supp} (\mathbf {N}) \cup \operatorname{Supp} (\mathbf {M} / \mathbf {N}).
\]

\begin{enumerate}
\def\labelenumi{(\roman{enumi})}
\setcounter{enumi}{1}
\tightlist
\item
  If \(\mathbf{M}\) is the sum of a family \((\mathbf{N}_{\iota})_{\iota \in \mathbf{I}}\) of submodules, then
\end{enumerate}

\[
\operatorname{Supp} (\mathbf {M}) = \bigcup_ {i \in I} \operatorname{Supp} \left(\mathrm{N} _ {i}\right).
\]

(i) From the exact sequence \(0 \to \mathbf{N} \to \mathbf{M} \to \mathbf{M} / \mathbf{N} \to 0\) , we derive, for every prime ideal \(\mathfrak{p}\) of \(\mathbf{A}\) , the exact sequence

\[
0 \rightarrow \mathrm{N} _ {p} \rightarrow \mathrm{M} _ {p} \rightarrow (\mathrm{M} / \mathrm{N}) _ {p} \rightarrow 0
\]

(§ 2, no. 4, Theorem 1). For \(\mathbf{M}_{\mathfrak{p}}\) to be reduced to 0, it is necessary and sufficient that \(\mathrm{N}_{\mathfrak{p}}\) and \((\mathrm{M} / \mathrm{N})_{\mathfrak{p}}\) be so. In other words, the relation \(\mathfrak{p} \notin \operatorname{Supp}(\mathbf{M})\) is equivalent to `` \(\mathfrak{p} \notin \operatorname{Supp}(\mathbf{N})\) and \(\mathfrak{p} \notin \operatorname{Supp}(\mathbf{M} / \mathbf{N})\) '', which proves (i).

(ii) For every prime ideal p of A, \(M_{p}\) is the sum of the family of submodules \((\mathbf{N}_{\iota})_{\mathfrak{p}}\) (§2, no. 4). To say that \(M_{p} \neq 0\) means that there exists \(\iota \in I\) such that \((\mathbf{N}_{\iota})_{\mathfrak{p}} \neq 0\) , whence (ii).

COROLLARY. Let A be a ring, M an A-module, \((m_{i})_{i\in I}\) a system of generators of M and \(a_{i}\) the annihilator of \(m_{i}\) . Then Supp(M) = \(\bigcup_{i\in I} V(a_{i})\) .

\(\operatorname{Supp}(\mathbf{M}) = \bigcup_{i \in I} \operatorname{Supp}(Am_i)\) by Proposition 16 (ii). On the other hand, \(Am_i\) is isomorphic to the A-module \(A/a_i\) and we have seen that

\[
\operatorname{Supp} (\mathrm{A} / \mathfrak {a} _ {\iota}) = \mathrm{V} (\mathfrak {a} _ {\iota})
\]

(Example above).

PROPOSITION 17. Let A be a ring, M an A-module and a its annihilator; if M is finitely generated, then \(\operatorname{Supp}(\mathbf{M}) = \mathrm{V}(\mathfrak{a})\) and \(\operatorname{Supp}(\mathbf{M})\) is therefore closed in \(\operatorname{Spec}(\mathbf{A})\) . Let \((m_i)_{1 \leqslant i \leqslant n}\) be a system of generators of M and let \(\mathfrak{a}_i\) be the annihilator of \(m_i\) ; then \(\mathfrak{a} = \bigcap_{i=1}^{n} \mathfrak{a}_i\) , hence \(\mathrm{V}(\mathfrak{a}) = \bigcup_{i=1}^{n} \mathrm{V}(\mathfrak{a}_i)\) (no. 3, equation (3)) and the proposition follows from the Corollary to Proposition 16.

COROLLARY 1. Let A be a ring, M a finitely generated A-module and a an element of A. For a to belong to every prime ideal of the support of M, it is necessary and sufficient that the homothety of M with ratio a be nilpotent.

It follows from Proposition 17 that the intersection of the prime ideals belonging to Supp(M) is the radical of the annihilator a of M (no. 3, Proposition 11 (ii)). To say that a belongs to this radical is equivalent to saying that there exist a power \(a^{k} \in a\) and hence that \(a^{k}M = 0\) .

COROLLARY 2. Let A be a Noetherian ring, M a finitely generated A-module and a an ideal of A. For \(\operatorname{Supp}(\mathbf{M}) \subset \mathrm{V}(\mathfrak{a})\) , it is necessary and sufficient that there exist an integer \(k\) such that \(\mathfrak{a}^k\mathbf{M} = 0\) .

If b is the annihilator of M, the relation \(\operatorname{Supp}(\mathbf{M}) \subset \mathbf{V}(\mathfrak{a})\) is equivalent to \(\mathbf{V}(\mathfrak{b}) \subset \mathbf{V}(\mathfrak{a})\) by Proposition 17 and hence to \(\mathfrak{a} \subset \mathfrak{r}(\mathfrak{b})\) , where \(\mathfrak{r}(\mathfrak{b})\) is the radical of b (no. 3, Corollary 2 to Proposition 11). Since A is Noetherian, this condition is also equivalent to the existence of an integer \(k > 0\) such that \(a^{k} \subset b\) (§ 2, no. 6, Proposition 15).

PROPOSITION 18. Let M, M' be two finitely generated modules over a ring A; then

\[
\operatorname{Supp} (\mathbf {M} \otimes_ {\mathbf {A}} \mathbf {M} ^ {\prime}) = \operatorname{Supp} (\mathbf {M}) \cap \operatorname{Supp} (\mathbf {M} ^ {\prime}).\tag{10}
\]

We need to prove that, if p is a prime ideal of A, the relations \((\mathbf{M} \otimes_{\mathbf{A}} \mathbf{M}')_{p} \neq 0\) and `` \(M_{p} \neq 0\) and \(M'_{p} \neq 0\) '' are equivalent. As the \(A_{p}\) -modules \(M_{p} \otimes_{A_{p}} M'_{p}\) and \((\mathbf{M} \otimes_{\mathbf{A}} \mathbf{M}')_{p}\) are isomorphic (§2, no. 7, Proposition 18), our assertion follows from the following lemma:

LEMMA 3. Let B be a local ring and E and E' two finitely generated B-modules. If E ≠ 0 and E' ≠ 0, then \(E \otimes_{B} E' \neq 0\).

Let \(k\) be the residue field of B. By virtue of § 3, no. 2, Proposition 4, \(k \otimes_{\mathbf{B}} \mathbf{E} \neq 0\) and \(k \otimes_{\mathbf{B}} \mathbf{E}' \neq 0\) ; then we deduce that

\[
(k \otimes_ {\mathbf {B}} \mathbf {E}) \otimes_ {k} (k \otimes_ {\mathbf {B}} \mathbf {E} ^ {\prime}) \neq 0
\]

(Algebra, Chapter II, § 3, no. 7). But, since the tensor product is associative (loc. cit., § 3, no. 8), this tensor product is isomorphic to

\[
E \otimes_ {\mathbf {B}} ((k \otimes_ {k} k) \otimes_ {\mathbf {B}} \mathbf {E} ^ {\prime}) = \mathbf {E} \otimes_ {\mathbf {B}} (k \otimes_ {\mathbf {B}} \mathbf {E} ^ {\prime})
\]

and therefore to \(k \otimes_{\mathbf{B}} (\mathbf{E} \otimes_{\mathbf{B}} \mathbf{E}')\) , whence the lemma.

COROLLARY. Let M be a finitely generated A-module and n its annihilator. For every ideal a of A, Supp(M/aM) = V(a) ∩ V(n) = V(a + n).

\(\mathbf{M} / a\mathbf{M} = \mathbf{M}\otimes_{\mathbf{A}}(\mathbf{A} / a)\) and \(\mathbf{A} / a\) is finitely generated.

PROPOSITION 19. Let A, B be two rings, \(\phi: \mathbf{A} \to \mathbf{B}\) a homomorphism and

\[
^ {a} \phi : \operatorname{Spec} (\mathbf {B}) \rightarrow \operatorname{Spec} (\mathbf {A})
\]

the continuous mapping associated with \(\phi\) (Proposition 13). For every A-module M, \(\operatorname{Supp}(\mathbf{M}_{(\mathbf{B})}) \subset {}^a\overline{\phi}^{-1}(\operatorname{Supp}(\mathbf{M}))\) ; if also M is finitely generated, then

\[
\operatorname{Supp} (\mathbf {M} _ {(\mathbf {B})}) = ^ {a} \bar {\phi} ^ {- 1} (\operatorname{Supp} (\mathbf {M})).
\]

Let \(q\) be a prime ideal of \(B\) and \(p = \overline{\phi}^{-1}(q)\) . Suppose that \(q\) belongs to \(\operatorname{Supp}(M_{(B)})\) ; then \(M_{(B)} \otimes_B B_q = (M \otimes_A B) \otimes_B B_q = M \otimes_A B_q = (M \otimes_A A_p) \otimes_A B_q\) , since the homomorphism \(A \to B \to B_q\) factors into \(A \to A_p \to B_q\) ( \(\S 2\) , no. 1, Proposition 2); the hypothesis \(M_{(B)} \otimes_B B_q \neq 0\) implies therefore \(M \otimes_A A_p \neq 0\) , whence the first assertion. As the homomorphism \(\phi_q: A_p \to B_q\) is local, the second assertion follows from the following lemma:

LEMMA 4. Let A, B be two local rings, \(\rho: \mathbf{A} \to \mathbf{B}\) a local homomorphism and E a finitely generated A-module. If \(\mathbf{E} \neq 0\) , then \(\mathrm{E}_{(\mathbf{B})} \neq 0\) .

Let m be the maximal ideal of A and k = A/m the residue field; the hypothesis implies that \(B \otimes_{A} k = B/mB \neq 0\) ; since the tensor product is associative, \((\mathbf{E} \otimes_{\mathbf{A}} \mathbf{B}) \otimes_{\mathbf{A}} k\) is isomorphic to \(\mathbf{E} \otimes_{\mathbf{A}} (\mathbf{B} \otimes_{\mathbf{A}} k)\) , hence also to \(\mathbf{E} \otimes_{\mathbf{A}} (k \otimes_{k} (\mathbf{B} \otimes_{\mathbf{A}} k))\) and finally to \((\mathbf{E} \otimes_{\mathbf{A}} k) \otimes_{k} (\mathbf{B} \otimes_{\mathbf{A}} k)\) ; by §3, no. 2, Proposition 4, \(E \otimes_{A} k \neq 0\) , hence \((\mathbf{E} \otimes_{\mathbf{A}} \mathbf{B}) \otimes_{\mathbf{A}} k \neq 0\) (Algebra, Chapter II, §3, no. 7) and a fortiori \(E \otimes_{A} B \neq 0\) .

PROPOSITION 20. Let A be a ring and M a finitely generated A-module. For every prime ideal \(\mathfrak{p} \in \operatorname{Supp}(\mathbf{M})\) , there exists a non-zero A-homomorphism \(w: \mathbf{M} \to \mathbf{A}/\mathfrak{p}\) .

Let \(\mathfrak{p} \in \operatorname{Supp}(\mathbf{M})\) . As \(\mathbf{M}\) is finitely generated and \(\mathbf{M}_{\mathfrak{p}} \neq 0\) ,

\[
\mathbf {M} _ {\mathfrak {p}} / \mathfrak {p} \mathbf {M} _ {\mathfrak {p}} = \mathbf {M} _ {\mathfrak {p}} \otimes_ {\mathbf {A}} (\mathbf {A} _ {\mathfrak {p}} / \mathfrak {p} \mathbf {A} _ {\mathfrak {p}}) \neq 0
\]

(§ 3, no. 2, Proposition 4). Let \(\mathbf{K} = \mathbf{A}_{\mathfrak{p}} / \mathfrak{p}\mathbf{A}_{\mathfrak{p}}\) be the field of fractions of the integral domain \(\mathbf{A} / \mathfrak{p}\) ; since \(\mathbf{M}_{\mathfrak{p}} / \mathfrak{p}\mathbf{M}_{\mathfrak{p}}\) is a vector space over \(\mathbf{K}\) , which is not reduced to 0, there exists a non-zero linear form \(u: \mathbf{M}_{\mathfrak{p}} / \mathfrak{p}\mathbf{M}_{\mathfrak{p}} \to \mathbf{K}\) . If \((x_{i})_{1 \leqslant i \leqslant n}\) is a system of generators of \(\mathbf{M}, \bar{x}_{i}\) the image of \(x_{i}\) in the \((\mathbf{A} / \mathfrak{p})\) -module \(\mathbf{M}_{\mathfrak{p}} / \mathfrak{p}\mathbf{M}_{\mathfrak{p}}\) , there exists an element \(\alpha \neq 0\) of \(\mathbf{A} / \mathfrak{p}\) such that the \(\alpha u(\bar{x}_i)\) belong to \(\mathbf{A} / \mathfrak{p}\) for \(1 \leqslant i \leqslant n\) ; hence \(v = \alpha u\) is a non-zero \((\mathbf{A} / \mathfrak{p})\) -linear mapping from \(\mathbf{M}_{\mathfrak{p}} / \mathfrak{p}\mathbf{M}_{\mathfrak{p}}\) to \(\mathbf{A} / \mathfrak{p}\) . The composition

\[
w: \mathrm{M} \longrightarrow \mathrm{M} _ {p} \longrightarrow \mathrm{M} _ {p} / p \mathrm{M} _ {p} \stackrel {v} {\longrightarrow} \mathrm{A} / p
\]

is therefore the required homomorphism.

